\section{Linear harmonic response of a cantilever beam}

This example excites the tip of a cantilever beam with a harmonic force and evaluates 201 frequencies between 0 and 500. Rayleigh damping regularizes the resonance peaks. The response is written as separate real and imaginary displacement, stress, and strain fields for every frequency.

\begin{exampledeck}
*MODEL, NAME=HARMONIC_BEAM
*NODE, NSET=NALL
1, 0, 0, 0
2, 250, 0, 0
3, 500, 0, 0
4, 750, 0, 0
5, 1000, 0, 0
*ELEMENT, TYPE=B33, ELSET=BEAM
1, 1, 2
2, 2, 3
3, 3, 4
4, 4, 5
*NSET, NAME=ROOT
1
*NSET, NAME=TIP
5
*MATERIAL, NAME=STEEL
*ELASTIC, TYPE=ISOTROPIC
210000.0, 0.3
*DENSITY
7.85e-9
*PROFILE, NAME=RECT_20x10
200.0, 6666.7, 1666.7, 8333.3
*BEAMSECTION, ELSET=BEAM, MATERIAL=STEEL, PROFILE=RECT_20x10
0.0, 0.0, 1.0
*SUPPORT, SUPPORT_COLLECTOR=BC
ROOT, 0.0, 0.0, 0.0, 0.0, 0.0, 0.0
*CLOAD, LOAD_COLLECTOR=EXCITATION
TIP, 0.0, 0.0, -100.0, 0.0, 0.0, 0.0
*LOADCASE, TYPE=LINEARHARMONIC, NAME=FREQUENCY_RESPONSE
*SUPPORTS
BC
*LOADS
EXCITATION
*SOLVER, DEVICE=CPU, METHOD=DIRECT
*CONSTRAINTMETHOD, TYPE=NULLSPACE
*FREQUENCIES, SCALE=LINEAR
0.0, 500.0, 201
*DAMPING, TYPE=RAYLEIGH
0.0, 1.0e-5
*END
\end{exampledeck}

\section{Dual-mortar contact between two C3D8 blocks}

This regression model isolates the current surface-to-surface contact formulation. The lower block is fixed. The upper block starts with a 0.5 gap, is guided laterally, and receives a prescribed downward displacement at its top face. Contact begins when the initial gap has closed; before that point the upper block can translate without developing a normal contact reaction.

\begin{exampledeck}
*MODEL, NAME=CONTACT_MORTAR_PARALLEL_BLOCKS_C3D8
*NODE, NSET=NALL
1,  0.0, 0.0,  -1.0
2,  1.0, 0.0,  -1.0
3,  1.0, 1.0,  -1.0
4,  0.0, 1.0,  -1.0
5,  0.0, 0.0,   0.0
6,  1.0, 0.0,   0.0
7,  1.0, 1.0,   0.0
8,  0.0, 1.0,   0.0
9,  0.0, 0.0,   0.5
10, 1.0, 0.0,   0.5
11, 1.0, 1.0,   0.5
12, 0.0, 1.0,   0.5
13, 0.0, 0.0, 100.5
14, 1.0, 0.0, 100.5
15, 1.0, 1.0, 100.5
16, 0.0, 1.0, 100.5

*ELEMENT, TYPE=C3D8, ELSET=LOWER_BLOCK
1, 1, 2, 3, 4, 5, 6, 7, 8
*ELEMENT, TYPE=C3D8, ELSET=UPPER_BLOCK
2, 9, 10, 11, 12, 13, 14, 15, 16

*NSET, NAME=LOWER_ALL
1, 2, 3, 4, 5, 6, 7, 8
*NSET, NAME=UPPER_ALL
9, 10, 11, 12, 13, 14, 15, 16
*NSET, NAME=UPPER_TOP
13, 14, 15, 16

*SURFACE, NAME=MASTER_FACE
1, 1, S2
*SURFACE, NAME=SLAVE_FACE
2, 2, S1

*MATERIAL, NAME=LINEAR
*ELASTIC, TYPE=ISOTROPIC
1000.0, 0.0
*SOLIDSECTION, ELSET=LOWER_BLOCK, MATERIAL=LINEAR
*SOLIDSECTION, ELSET=UPPER_BLOCK, MATERIAL=LINEAR

*CONTACT, MASTER=MASTER_FACE, SLAVE=SLAVE_FACE, PENALTY=1.0e5

*SUPPORT, SUPPORT_COLLECTOR=BC
LOWER_ALL, 0.0, 0.0, 0.0, NAN, NAN, NAN
UPPER_ALL, 0.0, 0.0, NAN, NAN, NAN, NAN
UPPER_TOP, NAN, NAN, -1.0, NAN, NAN, NAN

*LOADCASE, TYPE=NONLINEARSTATIC, NAME=MORTAR_BENCHMARK
*SUPPORTS
BC
*SOLVER, DEVICE=CPU, METHOD=DIRECT
*CONSTRAINTMETHOD, TYPE=NULLSPACE
*NONLINEAR, CONTROL=LOAD, INITIAL_INCREMENT=0.1, MINIMUM_INCREMENT=0.1, MAXIMUM_INCREMENT=0.1, MAX_INCREMENTS=10, ADAPTIVE=OFF, MAXITER=30, TOL=1e-10, REGULARIZE_ZERO_ROWS=ON
*END
\end{exampledeck}

For this one-face/one-face benchmark the projected overlap area is 1, contact must remain inactive while the gap is positive, and total master and slave contact resultants must be equal and opposite. The repository additionally contains \file{examples/10_contact/103_contact_mortar_parallel_blocks_c3d8_expected.md} with the analytic pre-contact and penalty-response checks used for regression testing.
